\chapter{Priorisierung: Gewichtete Rankings und PageRank}
 \label{chap:priorisierung}

\begin{myremark}[Algorithmusart in diesem Kapitel]
Funktionsart: \textbf{Priorisierung}. \\
Umsetzungsart: \textbf{regelbasiert und datengetrieben}.
\end{myremark}

\begin{myremark}
Dieses Kapitel behandelt Algorithmen, die nicht primär Kategorien
oder Zahlen vorhersagen, sondern Elemente in eine \textbf{Reihenfolge} bringen.
\end{myremark}

\section{Grundidee: Ordnen statt klassifizieren}

Bei einer Priorisierung wird nicht gefragt \enquote{Welche Klasse?} oder \enquote{Welcher Wert?},
sondern \enquote{Was soll zuerst kommen?}. Das Ergebnis ist eine \textbf{Rangliste}.

Typische Beispiele sind Suchmaschinen, Newsfeeds, Lernplattformen oder Support-Tickets.
Die Qualität eines Priorisierungsalgorithmus hängt davon ab, wie sinnvoll die Reihenfolge ist.

\section{Gewichtete Priorisierung}

\begin{mydefinition}[Gewichtete Priorisierung]
Bei einer gewichteten Priorisierung bekommt jedes Kriterium ein \textbf{Gewicht}.
Aus den Teilwerten $z_1, z_2, \ldots, z_n$ und Gewichten $w_1, w_2, \ldots, w_n$ entsteht ein Score:
\[
\text{Score}(x) = \sum_{i=1}^{n} w_i z_i
\]
Oft gilt $\sum_i w_i = 1$ und alle Teilwerte sind auf $[0,1]$ normiert.
\end{mydefinition}

\begin{mycaution}
Wenn Merkmale unterschiedliche Skalen haben, verzerren grosse Zahlen die Rangliste.
Darum müssen die Werte vor dem Zusammenfassen oft normalisiert werden.
Für negative Kriterien wie \enquote{Aufwand} oder \enquote{Kosten} wird der Wert invertiert,
z.\,B. durch $1-z$.
\end{mycaution}

\begin{myexample}
Drei Quellen für ein Schulprojekt werden nach vier Kriterien bewertet:
Relevanz, Aktualität, Vertrauenswürdigkeit und geringer Aufwand.
Die Gewichte seien $0.40$, $0.25$, $0.25$ und $0.10$.
\end{myexample}

\begin{myexercise}[Gewichtete Scores berechnen]
Drei Lernquellen wurden bereits auf $[0,1]$ normiert:

\begin{center}
\begin{tabular}{lcccc}
\toprule
Quelle & Relevanz & Aktualität & Vertrauen & Aufwand \\
\midrule
A & 0.85 & 0.50 & 0.75 & 0.30 \\
B & 0.75 & 0.85 & 0.70 & 0.20 \\
C & 0.80 & 0.65 & 0.95 & 0.60 \\
\bottomrule
\end{tabular}
\end{center}

Berechnen Sie den Score mit
\[
\text{Score} = 0.40\cdot R + 0.25\cdot A + 0.25\cdot V + 0.10\cdot(1-\text{Aufwand})
\]
und ordnen Sie die drei Quellen.
\begin{myanswer}
A: $0.40\cdot0.85 + 0.25\cdot0.50 + 0.25\cdot0.75 + 0.10\cdot0.70 = 0.7225$ \\
B: $0.40\cdot0.75 + 0.25\cdot0.85 + 0.25\cdot0.70 + 0.10\cdot0.80 = 0.7675$ \\
C: $0.40\cdot0.80 + 0.25\cdot0.65 + 0.25\cdot0.95 + 0.10\cdot0.40 = 0.7600$ \\
Rangfolge: B, C, A.
\end{myanswer}
\end{myexercise}

\begin{myexercise}[Gewichte diskutieren]
Was verändert sich, wenn man Aktualität stärker gewichtet als Vertrauenswürdigkeit?
Nennen Sie ein Beispiel, in dem diese Entscheidung sinnvoll sein kann.
\end{myexercise}

\section{PageRank: Priorisierung im Netzwerk}

\begin{mydefinition}[PageRank]
PageRank priorisiert Webseiten nicht nur nach der Anzahl eingehender Links,
sondern auch nach der \textbf{Qualität} dieser Links.
Eine Seite ist wichtig, wenn wichtige Seiten auf sie verweisen.
\end{mydefinition}

Die Grundidee ist ein \enquote{zufälliger Surfer}:
Mit Wahrscheinlichkeit $d$ folgt er einem Link, mit Wahrscheinlichkeit $1-d$ springt er zufällig auf eine neue Seite.
Der Score einer Seite $p$ ist dann:
\[
PR(p) = \frac{1-d}{N} + d \sum_{q \to p} \frac{PR(q)}{\text{outdeg}(q)}
\]
mit $N$ Seiten und Dämpfungsfaktor $d$ (typisch: $0.85$).

\begin{mycaution}
PageRank ist nicht neutral: Linkfarmen (also Seiten, die sich gegenseitig verlinken), gekaufte Links oder starke Netzwerke können die Rangfolge verzerren.
Darum braucht man Missbrauchserkennung und zusätzliche Qualitätsmerkmale.
\end{mycaution}

\begin{myexercise}[PageRank intuitiv lesen]
Betrachten Sie das Netzwerk:
\begin{itemize}
    \item A verlinkt auf B und C
    \item B verlinkt auf C
    \item C verlinkt auf A
    \item D verlinkt auf C
\end{itemize}
Wenn alle Seiten mit gleichem PageRank starten und $d=0.85$ gilt:
\begin{enumerate}
    \item Welche Seite dürfte nach einem Rechenschritt am höchsten priorisiert sein?
    \item Warum ist ein Link von einer Seite mit vielen Ausgängen weniger wert?
\end{enumerate}
\begin{myanswer}
Nach einem Iterationsschritt ist C am höchsten priorisiert, weil A, B und D auf C verweisen.
Ein Link von einer Seite mit vielen Ausgängen wird auf viele Ziele verteilt und zählt deshalb weniger stark.
\end{myanswer}
\end{myexercise}

\begin{mychallenge}[Priorisierung vergleichen]
Vergleichen Sie gewichtete Priorisierung und PageRank:
\begin{enumerate}
    \item Welche Methode ist leichter erklärbar?
    \item Welche Methode ist anfälliger für Manipulation?
    \item Wo würden Sie welche Methode im Schulkontext einsetzen?
\end{enumerate}
\end{mychallenge}

\begin{myremark}[Notebook]
Praktische Umsetzung inkl. Übungen:
\href{\notebookbaseurl06_Priorisierung_PageRank.ipynb}{\texttt{06\_Priorisierung\_PageRank.ipynb}}
\end{myremark}
